\section{Discussion}

The main point is simple. Adult zebrafish telencephalon and optic tectum are not merely adjacent regions with somewhat different proteoform lists; they occupy sharply separated proteoform regimes whose dominant marker architecture follows known regional biology. The original study made that outcome plausible. The evidence assembled here makes it much harder to ignore. Once the named markers are rewritten as a matched-versus-spillover problem, the regional story is no longer a handful of appealing examples. The alignment fractions, odds-ratio interval, protein-collapsed robustness, intensity-weighted follow-up, and stress-test table all point in the same direction.

The biggest new factual gain over earlier versions of this note comes from the released Tel2 and Teo2 tables. They let us move from article-only overlap counts to exact accession+proteoform overlap, duplicate-incidence summaries, explicit marker-family membership tables, and a machine-readable evidence-traceability file. That extra layer matters because it exposes a small but real provenance wrinkle: the article reports 35 shared proteoforms, whereas the released exact-ID tables expose 29 exact matches. The new discrepancy diagnostic suggests that this is largely a representation issue rather than a biological contradiction. When bracketed PTM annotations and punctuation are stripped while accessions are retained, the shared count rises to 44, placing the published 35 between strict and relaxed matching rules. I therefore read the mismatch as a reminder that aggregate figure counts and released region tables are not always numerically identical, not as evidence that the present biological claim is an artifact of bookkeeping. The important point is that every public view still lands in the same regime: very low proteoform overlap and much higher protein overlap.

The supplementary workbook makes that provenance story cleaner as well. Its note sheet states that the public tables were identified in TopPIC at 1\% proteoform-spectrum-match FDR and 5\% proteoform-level FDR, with duplicate quantification performed in TopDiff, and its \texttt{Tel2} and \texttt{Teo2} sheets are exact copies of the standalone PRIDE tables used here. That does not prove the article's aggregate count of 35 shared proteoforms was produced under the same final aggregation logic, but it does make a simple table-side threshold mismatch less plausible than a representation or reporting difference.

The scope, however, remains narrow. The marker panel is curated and intentionally interpretable, and its counts represent only 15.20\% of the total regional count mass used here. The analysis therefore says less about the full latent structure of adult zebrafish brain proteomics than about one concrete point: this public pilot already supports a region-function proteoform axis that is sharper and more composition-aware than the original narrative presentation suggested. The technical replicate summaries remind us that this is still a pilot, not a large statistical atlas.

The new sensitivity checks sharpen that boundary in a useful way. Moving one to three matched counts per functional axis into spillover weakens the alignment effect, but it does not erase it. Likewise, excluding the three most obvious motor-associated families still leaves a strong optic-tectum-versus-telencephalon contrast in the remaining panel. The family-size-preserving randomization test is especially reassuring because it addresses the circularity objection in the units the paper actually uses: marker families. A null that preserves family sizes and the global telencephalon/optic-tectum label totals still centers around 88 matched counts, not 125, and none of 200{,}000 draws reaches the observed alignment. The duplicate-incidence calculations also push in the same direction: Chao2-style richness inflation lowers the exact-ID Jaccard if shared overlap is held fixed, and even the generous scaled-shared scenarios keep the overlap under 0.04. The first-order jackknife bound does the same. The new row-resampling and run-pair sensitivity files reach the same qualitative conclusion rather than rescuing a high-overlap story, and the prevalence-adjusted centered-Jaccard test shows that even the protein-collapsed overlap lies far below what the fixed margins would predict under independence. The minimal MNAR-style fill barely moves the weighted similarities. That means the current regional-separation claim survives modest curation perturbations, a plausible contamination objection, and more than one missingness model, even though it remains a public-table argument rather than a raw-feature validation study.

The orthogonal biological context also looks better after taking the review criticism seriously. The telencephalon side of the story is already well supported by behavioral and transcriptomic work on social orienting, learning, and memory \citep{stednitz2018social,reemst2023memory,pandey2023telencephalon}. On the optic-tectum side, the reticulon, myosin, actin, and troponin families fit a region that is visually driven, circuit-rich, and heavily involved in orienting transformations \citep{roeser2003tectum,orger2016visuomotor}. That does not prove every high-abundance motor-associated proteoform is a clean neuronal signal rather than a sampling or purity issue, but it makes the contamination-only explanation less attractive, especially once the non-motor panel still aligns strongly after exclusion and the released tables do not show a matching optic-tectum enrichment for the small glial/myelin sentinel screen. The new independent panel helps here as well. It is deliberately not a second copy of the source paper's own functional storytelling; it is a composition-oriented panel built from adult zebrafish myelin and adult optic-tectum radial-glia literature \citep{siems2021myelin,ito2010tectum}. The fact that it still separates the regions strongly suggests that the central result is broader than one hand-curated interpretive choice.

The PTM follow-up is also more informative than a simple ``there are many modified proteoforms'' statement, but it contributes by setting a boundary on current knowledge rather than by extending the positive claim. The raw matched-marker counts still show an acetylation asymmetry, and that asymmetry survives the limited Holm correction against the single exploratory non-acetyl bucket. Once the review forced us to model the public detectability proxies we actually have, however, the signal weakened substantially. The acetylated telencephalon-matched rows are heavier, longer, and more N-terminal than their optic-tectum counterparts, and the coarse covariate-adjusted logistic sensitivity no longer supports a confident region effect. Restricting to first-residue$\leq2$ rows removes the enrichment entirely. I therefore no longer read the acetylation contrast as a robust regional PTM program. The honest interpretation is narrower: the public tables contain a potentially interesting acetylation asymmetry, but the current release does not let us separate biology cleanly from proteoform geometry and coverage bias. That negative clarification is part of the contribution of the paper, because it prevents a fragile PTM hint from being mistaken for a settled biological program.

\subsection{Limits of inference}

This is not a raw-data reprocessing study, a new mass-spectrometry experiment, or a proteome-wide statistical atlas. It does not estimate region-wide significance for every detected proteoform, recover the full TopPIC workflow, or reinterpret unidentified mass shifts. Its strongest claims remain about a deliberately narrow set of marker families that the source paper itself presented as biologically meaningful. Even so, that is enough to support a useful statement about regional molecular organization, drawn from published counts, released region tables, and known regional function claims.

Several of the external review's requests still point to the same missing ingredient: richer feature-level detection histories. The released duplicate tables were enough for first-pass Chao2 and jackknife sensitivity, normalization checks, explicit marker-family traceability, and a limited PTM screen, but they still do not expose the full per-feature FDRs, localization confidence, charge-state resolution, or cross-region replicate incidence that a raw TopPIC export would provide \citep{choi2022toppic}. They also do not expose the higher-frequency incidence counts that would be required for iChao-style corrections. The current package therefore cannot adjudicate every possible redundancy or detectability objection. It can only show that several obvious objections fail to overturn the regional-function signal.

The literature on label-free normalization and missing-value handling is broader than what this tiny public design can support. EigenMS and later missingness-aware or imputation-oriented workflows were built for larger matrices and richer replication structures than two duplicate runs per region \citep{karpievitch2009eigenms,karpievitch2012missing,lazar2016missing}. For that reason this paper does not pretend to estimate a censoring model or an imputation pipeline from the released tables. Instead it treats blank cells as absent detections in the baseline analysis, reports several identifiable scaling schemes, and adds a minimal low-intensity-limit sensitivity so the reader can see how little the intensity-based overlaps move under a simple MNAR-style alternative.

The same ``do only what the public data can support'' rule also shapes the notation discussion. We now report that the relaxed overlap diagnostic is stable across three simple deterministic encodings, but that is still not the same thing as a full ProForma 2.0 implementation or a Proteoform Registry mapping \citep{leduc2018proforma}. Those standards are the right long-term answer when source strings and registry identifiers are available. Likewise, complementary bottom-up postquantification frameworks such as HiQuant are useful reference points for future cross-checking of proteoform stoichiometry, but they need peptide-level inputs that are outside the scope of this public top-down table release \citep{bryan2016hiquant}.

Those limits point directly to future work. The obvious next step is to extend the same analysis to the raw duplicate files now located on the PRIDE landing page together with the corresponding TopPIC feature exports or search intermediates, and then ask how much of the same function-legibility argument survives a proteome-wide treatment. Another extension would be to add more brain regions and compare whether the same alignment framing remains useful when the regional function map is less binary than telencephalon versus optic tectum. A third would be to formalize tissue-purity checks for the motor-associated families so that the contamination question can move from a stress test into a direct measurement. A fourth would be to connect this type of proteoform-specific auditability note more explicitly to the broader spatial-omics ecosystem: KRONOS and HEIST show what foundation-model scale looks like for spatial proteomics, imaging, and transcriptomics \citep{shaban2025kronos,madhu2025heist}, whereas this note occupies a complementary niche in which the priority is transparent proteoform-specific accounting over scale.
